% source: https://github.com/pfradin/luadraw/discussions/95#discussioncomment-14597160 (pfradin, block 1)
\documentclass[12pt]{standalone}
\usepackage[svgnames]{xcolor}
\usepackage[3d]{luadraw}
\begin{document}

\begin{luacode*}
function define_getcolor(F, pal, mode, default_color)
-- F = list of facets
-- pal = palette of colors (rgb tables)
-- mode = "x", "y", or "z"
-- returns the fonction : getcolor(facet)
    default_color = default_color or White
    local x1,x2,y1,y2,z1,z2 = getbounds3d(F)
    local getcolor = function(f) -- f is a facet, returns the color of f
        local A = isobar3d(f)
        if mode == "x" then return palette(pal,(A.x-x1)/(x2-x1),true)
        elseif mode == "y" then return palette(pal,(A.y-y1)/(y2-y1),true)
        elseif mode == "z" then return palette(pal,(A.z-z1)/(z2-z1),true)
        else return default_color
        end
    end
    return getcolor
end

function graph3d:adjust_color(F,color,contrast,twoside) -- F = facet, contrast in [0,1], twoside=true/false
-- adjust color based on facet normal vector (scalar product between the normal vector of the facet and the vector directed towards the observer.)
-- returns a color, normal vector and coef
    local coef, N, r, g, b = {}
    local A, B, C, k, n, m, coef, ok, newcolor
    k = 2; m = #F; ok = false
    while (not ok) and (k < m) do
        A, B, C = F[1], F[k], F[k+1]
        N = pt3d.normalize(pt3d.prod(B-A,C-A))
        coef = pt3d.dot(self.Normal,N)
        ok = (coef ~= nil)
        k = k+1
    end
    if ok then 
        if math.abs(coef) < 1e-8 then coef = 0 end
        local neg = twoside and (coef < 0)
        local c = round(math.exp( math.log(math.abs(coef))/2.5*contrast),4)
        if notDef(c) then c = 0 end
        if type(color) == "string" then -- color must be a color name
            if neg then
                newcolor = color.."!50!white!"..strReal(100*c).."!black"
            else
                newcolor = color.."!"..strReal(100*c).."!black"
            end
        else 
            r, g, b = table.unpack(color)
            if neg then 
                newcolor = rgb({c*(1+r)/2,c*(1+g)/2,c*(1+b)/2})
            else
                newcolor = rgb({c*r,coef*g,c*b})
            end
        end
        return newcolor, N, coef
    end
end

function graph3d:drawfacet(S,args) -- internal use, S is a list of already sorted facets
    local oldfillstyle = self.param.fillstyle
    local oldfillopacity = self.param.fillopacity
    local oldfillcolor = self.param.fillcolor
    local oldlinestyle = self.param.linestyle
    local oldlineopacity = self.param.lineopacity
    local oldlinecolor = self.param.linecolor
    local oldlinewidth = self.param.linewidth
    local coul = args.color
    local getcolor
    local usepalette = args.usepalette -- palette = {Pal, "x", or "y" or "z"}, Pal is a list of colors in RGB table format
    if usepalette ~= nil then
        local pal, mode = table.unpack(usepalette)
        getcolor = define_getcolor(S,pal,mode,coul)
    end
 
    if (args.mode == mShadedOnly) and (args.opacity ~= 1)  then self:Linestyle("noline") end
    if (args.mode ~= mShadedOnly) then self:Lineoptions(args.edgestyle,args.edgecolor,args.edgewidth) end
    
    -- dessin des facettes
    if (usepalette == nil) and ((args.mode == mFlat) or (args.mode == mFlatHidden)) then -- facets of the same color defined by args.color
        if type(coul) == "table" then coul = rgb(coul) end
        self:Filloptions("full",coul,args.opacity)
        self:Dpolyline3d(S,true)
    elseif usepalette ~= nil then  -- color of each facet chosen from a palette
        for _,F in ipairs(S) do
            if (args.mode == mFlat) or (args.mode == mFlatHidden) then
                coul = rgb(getcolor(F))
            else
                coul = self:adjust_color(F, getcolor(F),args.contrast,args.twoside)
            end
            self:Filloptions("full",coul,args.opacity)
            if (args.mode == 5) then 
                if (args.opacity == 1) then
                    self:Lineoptions("solid",coul,1); self:Lineopacity(args.opacity)
                    self:Dpolyline3d(F,true)
                else
                    self:Dpolyline3d(F,true)
                end
            else
                self:Dpolyline3d(F,true)
            end
        end
    else --modes 3, 4 ou 5 couleur des facettes uniques mais nuancées en fonction du vecteur normal
        for _,F in ipairs(S) do
            coul = self:adjust_color(F,args.color,args.contrast,args.twoside)
            if coul ~= nil then
                self:Filloptions("full",coul,args.opacity)
                if (args.mode == mShadedOnly) then 
                    if (args.opacity == 1) then
                        self:Lineoptions("solid",coul,1); self:Lineopacity(args.opacity)
                        self:Dpolyline3d(F,true)
                    else
                        self:Dpolyline3d(F,true)
                    end
                else
                    self:Dpolyline3d(F,true)
                end
            end
        end
    end
    self:Filloptions(oldfillstyle,oldfillcolor,oldfillopacity)
    self:Lineoptions(oldlinestyle,oldlinecolor,oldlinewidth); 
    self:Lineopacity(oldlineopacity)
end

function graph3d:Dmixfacet(...) --Dmixfacet(F1,args1, F2,args2, ...)
-- dessine une liste de facettes F (avec coordonnées 3d)
-- args est une table à 10 champs :
-- args = {mode=0/1/2/3/4/5, contrast=1, backcull=true/false, edgestyle=defaut, edgecolor=default, edgewidth=defaut, twoside=true/false, color=White, opacity=1, usepalette=nil}    
    --mode=0: arêtes seulement
    --mode=1 et 2: faces peintes couleur unie+arêtes visibles
    --mode=3 et 4 faces peintes couleur nuancée+arêtes visibles
    --mode=5 faces peintes couleur nuancée pas d'arêtes
    local mode = 3
    local contrast = 1
    local edgestyle = self.param.linestyle
    local edgecolor = self.param.linecolor
    local edgewidth = self.param.linewidth
    local backcull = false
    local twoside = true
    local color = "white"
    local opacity = 1
    local S, face, args = {}
    for k,F in ipairs{...} do
        if k%2 == 1 then 
            face = F
            if not isID3d(self.matrix3d) then face = self:Mtransform3d(face) end -- on applique la matrice de transformation
        else
            args = F
            if args.mode == nil then args.mode = mode else mode = args.mode end
            if args.contrast == nil then args.contrast = contrast else contrast = args.contrast end
            if args.edgestyle == nil then args.edgestyle = edgestyle else edgestyle = args.edgestyle end
            if args.edgecolor == nil then args.edgecolor = edgecolor else edgecolor = args.edgecolor end
            if args.edgewidth == nil then args.edgewidth = edgewidth else edgewidth = args.edgewidth end
            if args.opacity == nil then args.opacity = opacity else opacity = args.opacity end
            if args.backcull == nil then args.backcull = backcull else backcull = args.backcull end
            if args.color == nil then args.color = color else color = args.color end
            if args.twoside == nil then args.twoside = twoside else twoside = args.twoside end
            if args.usepalette == nil then
                args.getcolor = function(f)
                    return args.color
                end
            else
                local pal, mode = table.unpack(args.usepalette)
                args.getcolor = define_getcolor(face,pal,mode,args.color)
            end
            if isPoint3d(face[1]) then face = {face} end
            for _, f in ipairs(face) do
                table.insert(f,args) -- chaque facette est accompagnée de ses arguments
                table.insert(S,f)
            end
        end
    end
    local oldmatrix = self.matrix3d
    self.matrix3d = ID3d
    local oldfillstyle = self.param.fillstyle
    local oldfillopacity = self.param.fillopacity
    local oldfillcolor = self.param.fillcolor
    local oldlinestyle = self.param.linestyle
    local oldlineopacity = self.param.lineopacity
    local oldlinecolor = self.param.linecolor
    local oldlinewidth = self.param.linewidth  
    local coul
    S = self:Sortfacet(S)
    for _,F in ipairs(S) do
        args = table.remove(F)
        self:Lineoptions(args.edgestyle,args.edgecolor,args.edgewidth)
        -- dessin
        if (args.mode == mWireframe) then -- arêtes seulement (mWireframe)
            self:Filloptions("none")
            self:Dpolyline3d(F,true)
        elseif (args.mode == mFlat) or (args.mode == mFlatHidden) then --faces unies
            coul = args.getcolor(F)
            if type(coul) == "table" then coul = rgb(coul) end
            self:Filloptions("full",coul,args.opacity)
            self:Dpolyline3d(F,true)
        else --mode = 3,4 ou 5 faces peintes
            coul = self:adjust_color(F,args.getcolor(F),args.contrast,args.twoside)
            if coul ~= nil then
                self:Filloptions("full",coul,args.opacity)
                if (not args.backcull) or (not neg) then 
                    if args.mode == 5 then 
                        if (args.opacity == 1) then
                            self:Lineoptions("solid",coul,1); self:Lineopacity(args.opacity)
                            self:Dpolyline3d(F,true)
                        else
                            self:Linestyle("noline")
                            self:Dpolyline3d(F,true)
                        end   
                    else
                        self:Dpolyline3d(F,true)
                    end
                end
            end
        end
    end
    self:Filloptions(oldfillstyle,oldfillcolor,oldfillopacity)
    self:Lineoptions(oldlinestyle,oldlinecolor,oldlinewidth); self:Lineopacity(oldlineopacity)
    self.matrix3d = oldmatrix
end

function graph3d:addFacet(facet,args)
-- ajouter des facettes à la scène 3d
-- facet : une facette ou une liste de facettes (avec points 3d)
-- args table à 12 champs
-- {color="white", opacity=1, boxed=false, backcull=false, edgewidth=6, edge=false, edgecolor="black", hidden=false, hiddenstyle="dotted", contrast=1, twoside=false, matrix=ID3d, usepalette=nil}
    if facet == nil then return {{"nul"}} end
    args = args or {}
    if isPoint3d(facet[1]) then facet = {facet} end
    local color = args.color or "White"
    local opacity = args.opacity or 1
    local backcull = args.backcull or false
    local boxed = args.boxed or false
    local contrast = args.contrast or 1
    local twoside = args.twoside
    local hidden = args.hidden
    if hidden == nil then hidden = Hiddenlines end
    local hiddenstyle = args.hiddenstyle or Hiddenlinestyle
    if twoside == nil then twoside = true end
    local edge = args.edge or false
    local edgecolor = args.edgecolor or self.param.linecolor
    local edgewidth = args.edgewidth or self.param.linewidth
    local matrix = args.matrix or ID3d
    matrix = composematrix3d(self.matrix3d,matrix)
    local oldmatrix = self.matrix3d
    self.matrix3d = ID3d
    local F = facet
    if not isID3d(matrix) then 
        F = mtransform3d(F,matrix)
    end
    local x1,x2,y1,y2,z1,z2, getcolor 
    if args.usepalette == nil then
        getcolor = function(f)
            return color
        end
    else
        local pal, mode = table.unpack(args.usepalette)
        getcolor = define_getcolor(F,pal,mode,color)
    end    
    local res = {}
    local rep = {"facet"}
    if boxed then 
        if x1 == nil then x1,x2,y1,y2,z1,z2 = getbounds3d(F) end
        local eps = 1e-4
        x1 = x1-eps; x2 = x2+eps
        y1 = y1-eps; y2 = y2+eps
        z1 = z1-eps; z2 = z2+eps
        local P = parallelep(M(x1,y1,z1),M(x2-x1,0,0),M(0,y2-y1,0),M(0,0,z2-z1))
        insert(res, self:addWall(poly2facet(P),{matrix=ID3d}))
    end
    for _, face in ipairs(F) do
        local coul,n,coef = self:adjust_color(face,getcolor(face),contrast,twoside)
        if coul ~= nil then
            if coef < 0 then n = -n  end -- n must be directed towards the observer
            if (not backcull) or (coef > 0)  then            
                table.insert(rep, {face,{face[1],n},coul,opacity})
            end
        end
    end
    if edge then
        -- calcul des arêtes
        insert(res, self:addPolyline(facetedges(F),{color=edgecolor, hidden=hidden, hiddenstyle=hiddenstyle, width=edgewidth, matrix=ID3d}))
    end
    table.insert(res,rep) -- rep = {"facet", liste de faces}
    self.matrix3d = oldmatrix
    return res 
end
\end{luacode*}


 \begin{luadraw}{name=usepalette}
local P,bbox = read_obj_file("../obj/nefertiti.obj")
local g = graph3d:new{window3d=bbox,window={-7,5,-8,5},viewdir={35,60}, margin={0,0,0,0}, size={10,10}, bg="LightGray"}
local pal1 = {Gray, SlateGray, LightSkyBlue, LightPink, Pink, LightSalmon, FireBrick}
local pal2 = {Purple, Indigo, Blue, Green, Yellow, Orange, Red}
local pal3 = {LightSkyBlue,CornflowerBlue,LightCoral,IndianRed,DarkSalmon,YellowGreen,MediumSeaGreen}
local pal4 = {Cyan,DodgerBlue,Magenta,Red,Orange,Gold,LimeGreen}
local pal5 = {LightCyan,LightSkyBlue,Plum,LightCoral,LightSalmon,Khaki,PaleGreen}
g:Dpoly( shift3d(scale3d(P,0.5),-4*g:ScreenX()+2*g:ScreenY()), {usepalette={pal5,"z"}, mode=mShadedOnly})
local s = surface(function(x,y) return M(x,y,x^2-y^2)  end,-2,2,-2,2)
g:Setviewdir(60,60)
g:Dfacet( shift3d(scale3d(s,0.75),2*g:ScreenX()-2*g:ScreenY()), {usepalette={pal4,"z"}, mode=mShaded})
g:Show() 
\end{luadraw}
\end{document}
